% =====================================================================
% TERMO DE CONSENTIMENTO LIVRE E ESCLARECIDO (TCLE)
% Projeto guarda-chuva -- Cozy Pong / UFAL
%
% COMO COMPILAR:  latexmk -pdf TCLE.tex     (depois: latexmk -c)
%
% ANTES DE SUBMETER AO CEP: preencher todos os campos marcados em
% vermelho com a macro \pre{...}. Localize-os com:
%     grep -n "pre{" TCLE.tex
% =====================================================================
\documentclass[12pt,a4paper]{article}

\usepackage[T1]{fontenc}
\usepackage[brazil]{babel}
\usepackage{lmodern}
\usepackage[margin=2.4cm,bottom=3.0cm,top=2.4cm]{geometry}
\usepackage{amssymb}
\usepackage{enumitem}
\usepackage{xcolor}
\usepackage{fancyhdr}
\usepackage{lastpage}
\usepackage{microtype}
\usepackage{parskip}
\usepackage[hidelinks]{hyperref}

% ---- campo a preencher: aparece em vermelho, impossivel de passar batido
\newcommand{\pre}[1]{\textcolor{red}{\textbf{[#1]}}}

% ---- linha de assinatura
\newcommand{\linha}[1]{\rule{#1}{0.4pt}}

% ---- caixa de marcacao
\newcommand{\cx}{$\square$~}

% ---- rodape com rubricas e paginacao (exigencia usual de CEP)
\pagestyle{fancy}
\fancyhf{}
\renewcommand{\headrulewidth}{0pt}
\renewcommand{\footrulewidth}{0.4pt}
\fancyfoot[L]{\footnotesize Rubrica do participante: \linha{2.2cm} \quad
              Rubrica do pesquisador: \linha{2.2cm}}
\fancyfoot[R]{\footnotesize Página \thepage\ de \pageref{LastPage}}

\setlist[itemize]{leftmargin=1.4em,itemsep=2pt,topsep=3pt}

\begin{document}

% =====================================================================
\begin{center}
{\large\bfseries TERMO DE CONSENTIMENTO LIVRE E ESCLARECIDO (TCLE)}\\[4pt]
{\footnotesize Universidade Federal de Alagoas --- Instituto de Computação}\\[2pt]
{\footnotesize Versão \pre{1.0} --- \pre{data da versão}}
\end{center}

\vspace{4pt}
\hrule
\vspace{10pt}

% =====================================================================
\section*{1. Identificação}

\noindent\textbf{Título da pesquisa:} \pre{inserir o título aprovado}

\noindent\textbf{Pesquisador responsável:} Erick de Lima Mascarenhas\\
Curso de Ciência da Computação --- Instituto de Computação --- Universidade
Federal de Alagoas\\
Telefone: \pre{telefone com DDD} \quad E-mail: \pre{e-mail}

\noindent\textbf{Orientador(a):} \pre{nome completo e titulação}\\
Telefone: \pre{telefone} \quad E-mail: \pre{e-mail}

\noindent\textbf{Instituição proponente:} Universidade Federal de Alagoas (UFAL)

\noindent\textbf{Local de realização:} \pre{nome e endereço do laboratório ou sala}

% =====================================================================
\section*{2. Convite}

Você está sendo convidado(a) a participar, como voluntário(a), de uma pesquisa
acadêmica. Este documento explica por que a pesquisa está sendo feita, o que
você fará se aceitar participar, quais são os riscos e os benefícios, e quais
são os seus direitos.

\textbf{Leia com calma.} Pergunte tudo o que quiser, quantas vezes quiser, antes
de decidir. Você pode levar este documento para casa, conversar com quem quiser
e responder depois. Sua decisão de participar ou não \textbf{não traz nenhum
prejuízo} para você, e não afeta de nenhuma forma sua vida acadêmica, suas notas
ou sua relação com a universidade.

Se aceitar, você assinará \textbf{duas vias} deste termo, ambas rubricadas em
todas as páginas. Uma via fica com você; a outra, com o pesquisador.

% =====================================================================
\section*{3. Por que esta pesquisa está sendo feita}

Estresse e ansiedade são queixas comuns entre estudantes universitários. Existem
várias formas de tentar relaxar depois de um momento tenso: ouvir música,
meditar, praticar um esporte, jogar. O que ainda não se sabe bem é \textbf{quanto
cada uma dessas atividades ajuda}, e se uma ajuda mais que a outra.

Esta pesquisa desenvolveu um jogo de realidade virtual, chamado \textit{Cozy
Pong}, que mistura tênis de mesa com jogo de ritmo e foi projetado para ser
relaxante. O objetivo é medir se jogá-lo ajuda uma pessoa a se recuperar de um
momento de tensão, e comparar esse efeito com o de outras atividades.

Para que a comparação faça sentido, a pesquisa mede tanto \textbf{o que você
relata sentir} quanto \textbf{sinais do seu corpo}, como os batimentos cardíacos.

% =====================================================================
\section*{4. O que você fará se aceitar participar}

\subsection*{4.1 Quantas vezes}

Você participará de \textbf{1 a 6 encontros}, dependendo do estudo específico
para o qual for convidado(a). Cada encontro dura entre \textbf{30 e 120
minutos}. Quando houver mais de um encontro, eles serão em dias diferentes, com
pelo menos 24 horas de intervalo, em horários combinados com você.

Você será informado(a), antes de aceitar, de quantos encontros são no seu caso.

\subsection*{4.2 As atividades}

Em cada encontro, você fará \textbf{uma ou mais} das atividades abaixo. Nem todo
participante faz todas: quais você fará será informado no convite e novamente no
dia. Todas são atividades comuns de lazer, relaxamento ou esporte recreativo.

\textbf{Atividades em realidade virtual} (usando óculos de realidade virtual,
sentado ou em pé, em espaço livre de obstáculos):

\begin{itemize}
  \item[\textbf{a)}] Jogar o \textit{Cozy Pong} em versão calma, com música tranquila, ritmo lento e sem possibilidade de perder.
  \item[\textbf{b)}] Jogar o \textit{Cozy Pong} em versão mais movimentada, com música mais rápida, ritmo mais intenso e possibilidade de encerrar a partida por erros.
  \item[\textbf{c)}] Participar de uma sessão de meditação ou relaxamento guiado por áudio, sentado, em ambiente virtual tranquilo.
  \item[\textbf{d)}] Assistir passivamente a um ambiente ou vídeo em realidade virtual, sem interação.
  \item[\textbf{e)}] Jogar outro jogo de realidade virtual, usado para comparação.
\end{itemize}

\textbf{Atividades sem realidade virtual:}

\begin{itemize}
  \item[\textbf{f)}] Ouvir música sentado(a), com ou sem fones, sem nenhuma tarefa.
  \item[\textbf{g)}] Participar de uma sessão de meditação ou relaxamento guiado por áudio, sentado, sem óculos.
  \item[\textbf{h)}] Permanecer em repouso sentado(a), sem tarefa, por um período determinado.
  \item[\textbf{i)}] Jogar uma partida ou peloteio de \textbf{tênis de mesa real}, com raquete e bola convencionais, contra o pesquisador ou contra outro participante, em ritmo recreativo.
  \item[\textbf{j)}] Realizar atividade física leve, como caminhada em ritmo confortável ou alongamento.
  \item[\textbf{k)}] Jogar um jogo digital em tela comum, usado para comparação.
\end{itemize}

\textbf{Tarefa de tensão (importante que você leia com atenção):}

\begin{itemize}
  \item[\textbf{l)}] Realizar uma \textbf{tarefa de cálculo mental com tempo curto}, por cerca de 5 minutos, no computador.
\end{itemize}

Sobre esta tarefa, queremos ser completamente transparentes: \textbf{ela é
difícil de propósito}. O tempo é curto de propósito. É esperado que praticamente
todas as pessoas errem boa parte das contas, independentemente de habilidade com
matemática. Isso não é uma avaliação da sua capacidade.

O objetivo dela é gerar, por alguns minutos, aquela sensação de pressão que se
sente no dia a dia, para que possamos medir como as atividades ajudam na
recuperação depois. \textbf{Seu desempenho não será analisado individualmente
nem mostrado a ninguém.} Ao final de cada encontro, isso será explicado
novamente.

Você pode pedir para interromper esta tarefa a qualquer momento.

\subsection*{4.3 O que será medido}

\textbf{Questionários.} Você responderá questionários curtos, em papel ou no
computador, sobre como está se sentindo naquele momento, sobre como foi a
atividade e sobre alguns aspectos gerais do seu bem-estar recente. Não há
resposta certa ou errada. \textbf{Você pode deixar de responder qualquer
pergunta que não queira responder}, sem precisar justificar.

\textbf{Sinais do corpo.} Durante os encontros, poderão ser usados:

\begin{itemize}
  \item uma \textbf{cinta elástica no tórax} com sensores, que registra os batimentos do coração;
  \item \textbf{eletrodos nos dedos ou no punho}, que registram pequenas variações da pele relacionadas ao estado de ativação do corpo;
  \item uma \textbf{cinta elástica no abdômen ou tórax} para registrar a respiração;
  \item um sensor de \textbf{temperatura da pele};
  \item medida de \textbf{pressão arterial} com aparelho de braço, antes e depois da atividade.
\end{itemize}

\textbf{Nenhum desses sensores emite corrente elétrica, radiação ou qualquer
estímulo no seu corpo. Todos apenas registram.} Nenhum deles fura, corta ou
penetra a pele.

\textbf{Registro da atividade.} O jogo registra automaticamente informações da
partida, como acertos e tempo de reação. Os movimentos dos controles também
podem ser registrados.

\textbf{Gravação.} Os encontros poderão ser \textbf{gravados em áudio e/ou
vídeo}, apenas para conferir se os procedimentos foram seguidos corretamente e
para registrar observações. Você pode participar \textbf{recusando a gravação}, e
isso não impede sua participação. Marque sua escolha ao final deste termo.

% =====================================================================
\section*{5. Riscos e desconfortos}

Toda pesquisa com seres humanos envolve algum risco. Os riscos aqui são
considerados \textbf{baixos}, semelhantes aos de jogar um videogame, ouvir
música ou praticar um esporte recreativo. Ainda assim, listamos todos os que
identificamos, inclusive os pouco prováveis.

\textbf{Relacionados à realidade virtual:}
\begin{itemize}
  \item Enjoo, tontura, náusea, dor de cabeça, suor frio ou desorientação (chamados de \textit{cybersickness}), geralmente leves e passageiros
  \item Cansaço visual
  \item Desconforto pelo peso ou ajuste do óculos
  \item Risco de tropeço, colisão com móveis ou queda durante o uso
  \item Em pessoas com epilepsia fotossensível, risco de crise desencadeada por estímulos visuais (por isso essa condição é critério de não participação)
\end{itemize}

\textbf{Relacionados à atividade física (jogo em pé, tênis de mesa, caminhada):}
\begin{itemize}
  \item Cansaço, falta de ar, aumento dos batimentos
  \item Dor muscular durante ou nos dias seguintes
  \item Torção, distensão, queda ou impacto com a bola
  \item Em pessoas com problemas cardíacos não diagnosticados, risco de intercorrência cardiovascular durante o esforço (por isso há triagem de saúde antes)
\end{itemize}

\textbf{Relacionados à tarefa de cálculo mental:}
\begin{itemize}
  \item Sensação de pressão, frustração, irritação, nervosismo ou ansiedade durante os minutos da tarefa
  \item Desconforto por não conseguir acompanhar o ritmo da tarefa
\end{itemize}

\textbf{Relacionados aos questionários:}
\begin{itemize}
  \item Constrangimento, tristeza ou desconforto ao refletir sobre o próprio estado emocional
  \item Cansaço por responder perguntas repetidas
\end{itemize}

\textbf{Relacionados aos sensores:}
\begin{itemize}
  \item Desconforto, coceira ou vermelhidão na pele em contato com a cinta ou os eletrodos
  \item Constrangimento ao colocar a cinta torácica
\end{itemize}

\textbf{Relacionados aos dados:}
\begin{itemize}
  \item Risco de quebra de sigilo, ou seja, de alguém não autorizado ter acesso às suas respostas. Este risco é considerado baixo pelas medidas descritas na Seção 8, mas não pode ser eliminado por completo.
\end{itemize}

% =====================================================================
\section*{6. Como esses riscos serão reduzidos}

\begin{itemize}
  \item \textbf{Triagem antes de participar.} Pessoas com condições que tornem a participação insegura não serão incluídas, e isso será conversado com você.
  \item \textbf{Ambiente preparado.} O espaço de uso da realidade virtual será livre de obstáculos, com o pesquisador \textbf{presente e por perto durante toda a atividade}, atento a sinais de desequilíbrio.
  \item \textbf{Interrupção imediata.} A atividade é interrompida na hora se você pedir, por qualquer motivo e \textbf{sem precisar dar explicação}, ou se o pesquisador perceber qualquer sinal de mal-estar.
  \item \textbf{Pausa e retomada.} Você pode pedir pausas quando quiser.
  \item \textbf{Aquecimento e espaço adequado} antes das atividades físicas, com água disponível.
  \item \textbf{Higienização} do óculos, das cintas e dos sensores entre participantes. Eletrodos descartáveis ou higienizados conforme protocolo.
  \item \textbf{Explicação ao final} de cada encontro sobre a tarefa de cálculo, deixando claro que a dificuldade era proposital e que o desempenho não é avaliado.
  \item \textbf{Direito de não responder} qualquer pergunta.
  \item \textbf{Nenhuma pergunta ou atividade} de caráter íntimo, sexual, religioso, político ou sobre uso de substâncias ilícitas.
\end{itemize}

% =====================================================================
\section*{7. Assistência, acompanhamento e indenização}

Esta seção descreve garantias que você tem por lei.

\textbf{Assistência imediata.} Se você sentir qualquer mal-estar, desconforto
físico ou emocional durante ou logo após a participação, \textbf{o pesquisador
prestará assistência imediata}, interrompendo a atividade, acompanhando você e,
se necessário, acionando atendimento de saúde. Você não será liberado(a)
enquanto não estiver bem.

\textbf{Assistência integral e gratuita.} Você tem direito a \textbf{assistência
integral e gratuita} por qualquer dano decorrente da sua participação, pelo tempo
que for necessário, \textbf{imediato ou tardio}. Isso inclui encaminhamento e
acompanhamento para atendimento médico, fisioterapêutico ou psicológico,
conforme o caso, sem qualquer custo para você.

\textbf{Como acionar.} Se perceber qualquer problema que você atribua à
participação, mesmo dias ou semanas depois, entre em contato pelos telefones da
Seção 11. O pesquisador responsável se compromete a atender e providenciar o
encaminhamento.

\textbf{Encaminhamento em saúde mental.} Se, ao longo da participação, você
manifestar ou relatar sofrimento psíquico significativo, ou se suas respostas aos
questionários indicarem essa possibilidade, você será informado(a)
reservadamente e receberá informação sobre serviços de apoio psicológico
disponíveis, incluindo os serviços da própria universidade e a rede pública de
saúde. \textbf{Isso não é diagnóstico}, e os questionários desta pesquisa não
servem para diagnosticar nada.

\textbf{Indenização.} Você tem direito a \textbf{buscar indenização} por
eventuais danos decorrentes da sua participação nesta pesquisa, conforme a
legislação vigente. A assinatura deste termo \textbf{não implica renúncia a
nenhum direito seu}, inclusive o de recorrer judicialmente.

\textbf{Ressarcimento de despesas.} Sua participação \textbf{não terá nenhum
custo para você}. Despesas que você tiver exclusivamente por causa da pesquisa,
como transporte até o local ou alimentação em razão do tempo de permanência,
serão \textbf{ressarcidas} mediante comprovação. Combine isso com o pesquisador
antes.

\textbf{Sem pagamento.} Você \textbf{não receberá pagamento} pela participação,
pois a legislação brasileira não permite remuneração de voluntários em pesquisa.
O ressarcimento acima não é pagamento: é apenas devolução de despesa.

% =====================================================================
\section*{8. Sigilo, privacidade e uso dos dados}

\textbf{Anonimização.} Seus dados serão identificados apenas por um
\textbf{código} (por exemplo, P01). Seu nome não aparecerá em nenhum arquivo de
dados, em nenhuma apresentação e em nenhuma publicação.

\textbf{Guarda.} O documento que liga seu nome ao código ficará guardado
separadamente dos dados, em local de acesso restrito ao pesquisador responsável e
ao orientador. Arquivos digitais ficarão em dispositivo protegido por senha.

\textbf{Divulgação.} Os resultados serão divulgados de forma \textbf{agregada},
ou seja, como médias e comparações entre grupos, nunca de forma individual. Nada
publicado permitirá identificar você.

\textbf{Compartilhamento.} Dados \textbf{anonimizados} poderão ser
compartilhados com outros pesquisadores ou disponibilizados publicamente em
repositórios científicos, prática que aumenta a transparência da ciência. Nesse
caso, o compartilhamento inclui apenas os dados já sem qualquer identificação.

\textbf{Guarda e descarte.} Os dados serão guardados por no mínimo
\textbf{5 anos} após o término da pesquisa e depois descartados de forma segura.
As gravações de áudio e vídeo, quando houver, serão apagadas após a análise, em
prazo não superior a \textbf{5 anos}.

\textbf{Usos futuros.} Os dados anonimizados poderão ser usados em análises
posteriores relacionadas ao mesmo tema, sempre no escopo desta pesquisa. Novos
estudos com objetivos diferentes exigirão novo consentimento.

% =====================================================================
\section*{9. Benefícios}

\textbf{Benefício direto.} Não há garantia de benefício direto para você. É
possível que você se sinta mais relaxado(a) após algumas das atividades, mas isso
é justamente o que a pesquisa investiga e não pode ser prometido.

\textbf{Benefícios indiretos.} Ao participar, você contribui para que se entenda
melhor se e como jogos e atividades de relaxamento ajudam pessoas a lidar com
momentos de tensão. Isso pode orientar o desenvolvimento de ferramentas de
bem-estar mais acessíveis no futuro.

Você também poderá conhecer uma tecnologia de realidade virtual e, se tiver
interesse, receberá os resultados gerais da pesquisa quando concluída.

% =====================================================================
\section*{10. Sua liberdade de decidir}

Sua participação é \textbf{totalmente voluntária}.

\begin{itemize}
  \item Você pode \textbf{recusar} participar, sem precisar justificar.
  \item Você pode \textbf{desistir a qualquer momento}, inclusive no meio de uma atividade, sem precisar justificar.
  \item Você pode pedir que \textbf{seus dados sejam excluídos} da pesquisa, mesmo depois de participar, desde que os resultados ainda não tenham sido publicados.
  \item Nada disso trará \textbf{qualquer prejuízo} para você: não afeta notas, avaliações, estágios, relação com professores ou qualquer aspecto da sua vida acadêmica.
\end{itemize}

Você pode fazer perguntas ao pesquisador antes, durante e depois da participação,
e tem direito a receber resposta.

% =====================================================================
\section*{11. Contatos}

\textbf{Em caso de dúvidas sobre a pesquisa, mal-estar ou para acionar
assistência:}

\noindent Erick de Lima Mascarenhas --- pesquisador responsável\\
Telefone: \pre{telefone} \quad E-mail: \pre{e-mail}

\noindent \pre{nome do(a) orientador(a)} --- orientador(a)\\
Telefone: \pre{telefone} \quad E-mail: \pre{e-mail}

\textbf{Em caso de dúvidas sobre seus direitos como participante, ou para
denunciar qualquer irregularidade:}

\noindent Comitê de Ética em Pesquisa da Universidade Federal de Alagoas
(CEP/UFAL)\\
Endereço: \pre{endereço completo}\\
Telefone: \pre{telefone} \quad E-mail: \pre{e-mail}\\
Horário de atendimento: \pre{horário}

\begin{quote}
\small O Comitê de Ética em Pesquisa é um colegiado independente, criado para
defender os interesses dos participantes de pesquisa, respeitando sua dignidade,
seus direitos e seu bem-estar. Você pode procurá-lo a qualquer momento, sem
passar pelo pesquisador.
\end{quote}

% =====================================================================
\section*{12. Consentimento}

Eu, abaixo assinado(a), declaro que:

\begin{itemize}
  \item li ou tive lido para mim este documento, por inteiro;
  \item tive oportunidade de fazer perguntas e todas foram respondidas de forma satisfatória;
  \item entendi em que consiste a pesquisa, o que farei, quais são os riscos e os benefícios;
  \item entendi que minha participação é voluntária e que posso desistir a qualquer momento sem prejuízo;
  \item entendi que não receberei pagamento e que terei despesas ressarcidas;
  \item entendi que tenho direito a assistência e a buscar indenização por eventuais danos;
  \item recebi uma via deste documento, assinada e rubricada.
\end{itemize}

\textbf{Concordo em participar desta pesquisa.}

\vspace{6pt}
\noindent Sobre a gravação de áudio e vídeo dos encontros:

\noindent\cx \textbf{Autorizo} a gravação \hspace{2cm} \cx \textbf{Não autorizo} a gravação

\vspace{6pt}
\noindent Sobre o compartilhamento dos meus dados anonimizados com outros
pesquisadores:

\noindent\cx \textbf{Autorizo} \hspace{2cm} \cx \textbf{Não autorizo}

\vspace{16pt}

\noindent Local e data: \linha{6cm}, \linha{1.2cm} / \linha{1.2cm} / \linha{2cm}

\vspace{18pt}

\noindent Nome do participante (legível): \linha{9cm}

\vspace{14pt}

\noindent Documento de identidade: \linha{6cm}

\vspace{22pt}

\noindent Assinatura do participante: \linha{9cm}

\vspace{26pt}

\noindent Nome do pesquisador: Erick de Lima Mascarenhas

\vspace{22pt}

\noindent Assinatura do pesquisador: \linha{9cm}

\vspace{16pt}

\begin{center}
\small\textit{Este documento é emitido em duas vias, ambas rubricadas em todas as
páginas. Uma via permanece com o participante e a outra com o pesquisador
responsável.}
\end{center}

\end{document}
